\phantomsection
\addcontentsline{toc}{section}{Lampiran A --- Repositori GitHub}
\label{app:github-repository}

\section*{A.1 Tautan Repositori GitHub}

Kode sumber alat \textit{pipeline} migrasi PHP yang dikembangkan dalam
penelitian ini tersedia secara publik di repositori GitHub berikut:

\begin{center}
  \url{https://github.com/MFarrelAkbar1/skripsi-php-migration}
\end{center}

Repositori ini memuat seluruh modul Python (\texttt{pipeline/}), konfigurasi
Rector dan PHPStan, \textit{ruleset} Semgrep lokal (\texttt{rules/}),
berkas contoh pengujian (\texttt{tests/}), serta dokumentasi pendukung
(\texttt{docs/}) yang digunakan selama proses penelitian.

% --------------------------------------------------------
\section*{A.2 Struktur Folder Proyek}
% --------------------------------------------------------

Berikut adalah struktur direktori lengkap dari repositori
\textit{pipeline} migrasi PHP.

\begin{lstlisting}[style=pseudocodestyle]
skripsi-php-migration/
|-- README.md                           <- Dokumentasi utama
|-- requirements.txt                    <- Dependensi Python
|-- composer.json                       <- Dependensi PHP (Rector, 
                                          PHPStan)
|
|-- pipeline/                           <- Semua modul Python
|   |-- main.py                         <- Entry point & orkestrator 
                                           6 step
|   |-- scanner.py                      <- Semgrep security scanning
|   |-- converter.py                    <- Rector PHP conversion
|   |-- analyzer.py                     <- PHPStan static analysis
|   |-- ai_engine.py                    <- Ollama / LLM integration
|   |-- iso_mapper.py                   <- Pemetaan ISO/IEC 27001:2022
|   |-- composer_analyzer.py            <- Analisis kompatibilitas 
                                            composer.json
|   |-- llm_quality_checker.py          <- Syntax validity rate 
                                           per model
|   |-- accuracy_checker.py             <- Metrik akurasi pipeline
|   |-- accuracy_validator.py           <- Validasi hasil migrasi
|   |-- capture_responses.py            <- Capture raw LLM responses
|   |-- llm_comparison_chart.py         <- Chart perbandingan LLM
|   |-- generate_condition_a_charts.py  <- Chart eksperimen kondisi A
|   |-- generate_condition_b_charts.py  <- Chart eksperimen kondisi B
|   `-- pipeline_result_chart.py        <- Chart hasil pipeline 
                                           keseluruhan
|
|-- input/                              <- PHP 7.4 asli 
                                           (tidak dimodifikasi)
|-- output/                           <- PHP 8.x hasil konversi Rector
|-- reports/                            <- JSON reports + PNG charts
|-- rules/                              <- Semgrep rules lokal (pinned)

|-- tests/
|   `-- sample_php74/                   <- Sample PHP untuk 
                                           pengujian cepat
|       |-- sample_login.php
|       |-- sample_realistic.php
|       `-- sample_upload.php
|
|-- docs/
|   |-- architecture.md
|   |-- iso_controls.md
|   `-- pseudocode_appendix.md
|
`-- vendor/                              <- Paket Composer
\end{lstlisting}

\bigskip

Keterangan singkat setiap direktori utama:

\begin{center}
\renewcommand{\arraystretch}{1.3}
\begin{tabular}{|p{3.5cm}|p{11cm}|}
\hline
\textbf{Direktori / Berkas} & \textbf{Keterangan} \\
\hline
\texttt{pipeline/} &
  Seluruh modul Python yang mengimplementasikan enam tahap pipeline:
  pra-pindai, konversi, pasca-pindai, analisis statis, rekomendasi AI,
  dan pemetaan ISO. \\
\hline
\texttt{input/} &
  Folder tempat file PHP 7.4 asli ditempatkan sebelum diproses.
  Isi folder ini \textbf{tidak pernah dimodifikasi} oleh pipeline. \\
\hline
\texttt{output/} &
  Hasil konversi Rector: file PHP 8.x yang sudah diperbarui sintaksnya. \\
\hline
\texttt{reports/} &
  Laporan JSON dan grafik PNG yang dihasilkan setelah pipeline selesai,
  termasuk laporan kepatuhan ISO dan perbandingan model LLM. \\
\hline
\texttt{rules/} &
  \textit{Ruleset} Semgrep yang telah di-\textit{pin} secara lokal agar
  hasil eksperimen dapat direproduksi tanpa bergantung pada versi
  \textit{ruleset} online. \\
\hline
\texttt{tests/
        sample\_php74/} &
  Tiga berkas PHP sampel untuk pengujian cepat tanpa harus menyiapkan
  seluruh dataset studi kasus. \\
\hline
\texttt{docs/} &
  Dokumentasi arsitektur sistem, daftar kontrol ISO yang dipetakan,
  dan pseudocode dalam format Markdown. \\
\hline
\texttt{vendor/} &
  Paket PHP yang diinstall oleh Composer (Rector, PHPStan). \\
\hline
\end{tabular}
\end{center}
